\documentclass[../libro.tex]{subfiles}

% This is a package that includes many logotypes about TeX.
\usepackage{hologo}

% This package makes long tables.
\usepackage{longtable}

\begin{document}

\ifSubfilesClassLoaded{\backmatter}{}

\chapter{Colophon: notes on typesetting}

\begingroup
\setlength{\parindent}{1.5em}%for 12pt

\providecommand*{\typesettingengine}{whatever}
\providecommand*{\authoroftypesettingengine}{whomever}
\ifxetex
    \renewcommand*{\typesettingengine}{\hologo{XeLaTeX}}
    \renewcommand*{\authoroftypesettingengine}{Jonathan Kew}
\fi
\ifluatex
    \renewcommand*{\typesettingengine}{\hologo{LuaLaTeX}}
    \renewcommand*{\authoroftypesettingengine}{Taco Hoekwater}
\fi

This book was typeset in \typesettingengine{},
written by \authoroftypesettingengine{}.
The logotype \typesettingengine{},
along with other ``special typographical'' logotypes,
was produced by the \hologo{LaTeX} package \textsf{hologo},
written by Heiko Oberdiek.
The \hologo{LaTeX} software used for this book
was written by Leslie Lamport.
The \hologo{TeX} software,
which forms the base for \hologo{LaTeX},
was written by Donald Knuth.

The main text fonts in this book are
the OpenType format versions of
XITS (designed by Khaled Hosny)
and Source Han Serif (designed by Adobe and Google).
The sans serif fonts in this book are
the OpenType format version of
Fira Sans (designed by Carrois Apostrophe)
and the TrueType format version of
Sarasa Mono (designed by Belleve).
The monospaced font in this book is
the TrueType format version of
Sarasa Mono (designed by Belleve).
% The mathematical fonts in this book are
% the OpenType format versions of
% XITS Math (designed by Khaled Hosny),
% TeX Gyre Termes Math (designed by \hologo{TeX} Users Group),
% and STIX Two Math (designed by STIX Pub).
The mathematical font in this book is
the OpenType format version of
XITS Math (designed by Khaled Hosny).
The \typesettingengine{} packages
\textsf{fontspec} and \textsf{unicode-math},
both written by Will Robertson,
were used to manage fonts.

The cover of this book was inspired by
that of \textit{The not so short introduction to \hologo{LaTeX}},
written by Tobias Oetiker and others.
The long tables in this book were produced
by the \hologo{LaTeX} package \textsf{longtable},
written by David Carlisle.
The figures in this book were produced
by the \hologo{LaTeX} package \textsf{Ti\textit{k}Z},
written by Till Tantau.
The colophon (which you are reading now)
was inspired by that of
\textit{Linear algebra done right},
written by Sheldon Axler.

For more information about typesetting of this book,
please check the source code of this book,
hosted at
\url{https://github.com/septsea/det}.

\endgroup

\clearpage
\thispagestyle{empty}

\vspace*{\fill}

\begin{longtable}[c]{lllll}
    \multicolumn{5}{c}{\textbf{The Greek alphabet}}
    \\[2ex]
    \multicolumn{2}{c}{letter}
          &
    \multicolumn{2}{c}{name}
          &
    \multicolumn{1}{c}{pronunciation}
    \\ \hline
    \endfirsthead

    \multicolumn{5}{r}{(continued from the previous page)}
    \\[2ex]
    \multicolumn{2}{c}{letter}
          &
    \multicolumn{2}{c}{name}
          &
    \multicolumn{1}{c}{pronunciation}
    \\ \hline
    \endhead

    Α     & α     & άλφα    & alpha   & \textit{AL-fuh}                   \\
    Β     & β     & βήτα    & beta    & \textit{BAY-tuh}                  \\
    Γ     & γ     & γάμμα   & gamma   & \textit{GAM-muh}                  \\
    Δ     & δ     & δέλτα   & delta   & \textit{DEL-tuh}                  \\
    Ε     & ε (ϵ) & έψιλον  & epsilon & \textit{EP-suh-lon}               \\
    Ζ     & ζ     & ζήτα    & zeta    & \textit{ZAY-tuh}                  \\
    Η     & η     & ήτα     & eta     & \textit{AY-tuh}                   \\
    Θ (ϴ) & θ (ϑ) & θήτα    & theta   & \textit{THAY-tuh}                 \\
    Ι     & ι     & ιώτα    & iota    & \textit{eye-OH-tuh}               \\
    Κ     & κ (ϰ) & κάππα   & kappa   & \textit{KAP-uh}                   \\
    Λ     & λ     & λάμδα   & lambda  & \textit{LAM-duh}                  \\
    Μ     & μ     & μι      & mu      & \textit{mew}                      \\
    Ν     & ν     & νι      & nu      & \textit{new}                      \\
    Ξ     & ξ     & ξι      & xi      & \textit{ksigh}                    \\
    Ο     & ο     & όμικρον & omicron & \textit{OM-uh-CRON}               \\
    Π     & π (ϖ) & πι      & pi      & \textit{pie}                      \\
    Ρ     & ρ (ϱ) & ρο      & rho     & \textit{row}                      \\
    Σ     & σ/ς   & σίγμα   & sigma   & \textit{SIG-muh}                  \\
    Τ     & τ     & ταυ     & tau     & \textit{tow} (as in \textit{cow}) \\
    Υ     & υ     & ύψιλον  & upsilon & \textit{OOP-suh-LON}              \\
    Φ     & φ (ϕ) & φι      & phi     & \textit{fi} (as in \textit{hi})   \\
    Χ     & χ     & χι      & chi     & \textit{ki} (as in \textit{hi})   \\
    Ψ     & ψ     & ψι      & psi     & \textit{psigh}                    \\
    Ω     & ω     & ωμέγα   & omega   & \textit{oh-MAY-guh}               \\
\end{longtable}

\vspace{.5\baselineskip}\mbox{}\vfill

\clearpage

\end{document}

In publishing, a colophon is a brief statement
containing information about the publication
of a book such as an "imprint".

As is mentioned above,
this colophon was inspired by the one in Axler's book.

Postscript: the Greek alphabet was added for reference.
